\documentclass[12pt]{article}
\usepackage{amsmath}
\usepackage{amsfonts}
\usepackage{amssymb}
\usepackage{geometry}
\usepackage{tikz}
\geometry{a4paper, margin=1in}
\pagenumbering{gobble}


\setlength\parindent{0pt}

\title{02YMECH - Homework 13}
\date{Assigned for the week of Dec 15, 2025}
\author{}
\begin{document}
\maketitle

\section*{Question}

Consider a rigid body with principal moments of inertia
\[
I_1<I_2<I_3,
\]
and body-fixed principal axes
\((\mathbf e_1,\mathbf e_2,\mathbf e_3)\).
The motion is torque-free.

At time \(t=0\), the components of the angular velocity in the body frame
are
\[
\omega_1(0)=\omega_0,
\qquad
\omega_2(0)=\omega_0,
\qquad
\omega_3(0)=0,
\]
with \(\omega_0>0\).

\begin{enumerate}
  \item[(a)] Write Euler’s equations for this system and show that both the
        kinetic energy \(T\) and the squared angular momentum \(L^2\) are
        conserved.

  \item[(b)] Using only these two conserved quantities, eliminate
        \(\omega_2\) and \(\omega_3\) to obtain a first-order differential
        equation for \(\omega_1(t)\) of the form
        \[
        \dot{\omega}_1^{\,2}=f(\omega_1),
        \]
        and determine the explicit polynomial \(f\).

  \item[(c)] Determine the range of allowed values of \(\omega_1(t)\) and show
        that the motion of \(\omega_1\) is periodic.

  \item[(d)] Express the period of \(\omega_1(t)\) as an integral,
        clearly stating the change of variables used. Do not evaluate the
        integral.
\end{enumerate}


\end{document}
